\section{Koszul duality and the inverse Varchenko matrix}\label{sec:koszul}

Let $(\M,0)$ be an affine oriented matroid. We use the definitions
\[
 V_b(P,Q)=b_{\Sep(P,Q)},\qquad
 W_q(P,Q)=(-q)^{d(P,Q)}h((\bar P\cap\bar Q)_\bd,q^2)
\]
for all $P,Q\in\cF$, interpreting faces as in
Section~\ref{sec:background}. We denote the submatrix of $W_q$ indexed by bounded topes by $S_q$. The variables $b_i$ in this section are indexed by $E_0$. The main goal of this section is to relate the matrices $V_b$ and $W_q$ in the case that $(\M,0)$ is normal crossing via categorification and Koszul duality.

\subsection{Taut path and intersection algebras}
Fix a field $k$. We compose paths from left to right, so $e_PAe_Q$
has initial region $P$ and terminal region $Q$.

\begin{definition}
The \emph{taut path algebra} is the vector space
\[
 \Taut=\bigoplus_{P,Q\in\cF}k\tau_{PQ}
\]
with multiplication
\[
 \tau_{PQ}\tau_{SR}=
 \begin{cases}
  \tau_{PR},&Q=S\text{ and }d(P,Q)+d(Q,R)=d(P,R),\\
  0,&\text{otherwise}.
 \end{cases}
\]
\end{definition}

For each $P$, the element $e_P=\tau_{PP}$ is idempotent, and
$\sum_{P\in\cF}e_P=1$. A \emph{path} is a sequence of regions
sharing successive facets. It is \emph{taut} if it crosses each
separating hyperplane once and no others. The element $\tau_{PQ}$
is represented by any taut path from $P$ to $Q$.

\begin{definition}
For $P,Q\in\cF$, let
\[
 R_{PQ}=k[x_i\mid i\in E_0]/(x_J\mid J\notin\Delta_{PQ}),
 \qquad x_J=\prod_{i\in J}x_i.
\]
Write $1_{PQ}$ for the class of $1$ in this block. We use
$x_\varnothing=1$, so $R_{PQ}=0$ when $P,Q$ have no common feasible
face. The \emph{chamber intersection algebra} is
\[
 \Inter=\bigoplus_{P,Q\in\cF}R_{PQ},
\]
where the variables act centrally, products with different middle
labels are zero, and
\begin{equation}\label{eq:intersection-product}
 1_{PQ}1_{QR}=x_{\Sep(P,Q)\cap\Sep(Q,R)}1_{PR}.
\end{equation}
\end{definition}
The products respect the defining ideals, since a common face of $P,R$
on the hyperplanes in $\Sep(P,Q)\cap\Sep(Q,R)$ is also a face of
$Q$. Moreover, one may show associativity by counting repeated crossings in each coordinate.

Again $e_P=1_{PP}$ is idempotent and $\sum_Pe_P=1$, so the two
algebras have the same degree-zero algebra $\bigoplus_Pke_P$.
For adjacent regions separated by $H_i$,
\[
 1_{PQ}1_{QP}=x_ie_P,\qquad \tau_{PQ}\tau_{QP}=0.
\]
\begin{remark}
For realizable normal crossing arrangements, $\Inter$ is the algebra
$\widetilde B_{\mathrm{ext}}$ of \cite[Section 4.5]{BLPWgale}.
Their $\widetilde B$ is its corner on the feasible chambers bounded
with respect to a chosen generic objective function. Imposing the
linear parameter relations then gives their $B$
\cite[Definition 4.3 and Remark 4.5]{BLPWgale}.
Kowalenko--Mautner use the same face rings and multiplication for
oriented matroid programs \cite[Definition 5.9]{KowalenkoMautner}.
Our algebra uses all feasible chambers and retains the parameters.
Boundedness for an objective function differs from the boundedness
defining $\cB$ here.
\end{remark}

Both algebras are graded by $\NN^{E_0}$, with
\[
 \mdeg\tau_{PQ}=\mdeg1_{PQ}=\one_{\Sep(P,Q)},\qquad
 \mdeg x_i=2\mathbf e_i.
\]
Here $\one_J$ is the indicator vector of $J$, and $\mathbf e_i$ is
the $i$th standard basis vector. Total degree is the sum of the
coordinates of the multidegree.

\subsection{Bases and Hilbert matrices}
The block $e_P\Taut e_Q$ has basis $\{\tau_{PQ}\}$, and
$e_P\Inter e_Q$ has basis
\[
 \{x^\nu1_{PQ}\mid\nu\in\NN^{E_0},\ \supp\nu\in\Delta_{PQ}\}
\]
since the defining monomial ideal of $R_{PQ}$ is spanned by precisely
the monomials whose supports do not belong to $\Delta_{PQ}$.
For either algebra $A$, define its multigraded Hilbert matrix by
\[
 \Hilb_b(A)_{P,Q}=\sum_{\mathbf m\in\NN^{E_0}}
       \dim_k(e_PA_{\mathbf m}e_Q)b^{\mathbf m},
 \qquad b^{\mathbf m}=\prod_{i\in E_0}b_i^{m_i}.
\]
The bases give
\begin{align*}
 \Hilb_b(\Taut)_{P,Q}&=b_{\Sep(P,Q)}=V_b(P,Q),\\
 \Hilb_b(\Inter)_{P,Q}&=b_{\Sep(P,Q)}
    \sum_{J\in\Delta_{PQ}}\prod_{i\in J}\frac{b_i^2}{1-b_i^2}.
\end{align*}
By applying Lemma~\ref{lem:face-enumeration} to this formula, we obtain the following relation with the $q$-intersection matrix.
\begin{corollary}\label{cor:graded-hilbert}
For total degree, $\Hilb_q(\Taut)=V_q$. If $(\M,0)$ is normal crossing, then
\[
 \Hilb_q(\Inter)=(1-q^2)^{-r}W_{-q}.
\]
\end{corollary}
\subsection{Inverse Varchenko matrix via Koszul duality}
Let $A$ be a nonnegatively graded $k$-algebra with finite-dimensional
graded pieces and $A_0=\bigoplus_{i\in I}ke_i$ for pairwise orthogonal
idempotents $e_i$. It is \emph{quadratic} if
$A=T_{A_0}(V)/(\mathcal R)$, where $V=A_1$ and
$\mathcal R\subseteq V\otimes_{A_0}V$. Its quadratic dual is
\[
 A^!=T_{A_0}(V^*)/(\mathcal R^\perp),
\]
where the bimodule dual reverses endpoints. Dual generators retain
their multidegrees \cite[Definition 2.8.1]{BeilinsonGinzburgSoergel1996}.
The algebra is \emph{Koszul} if $A_0$ has a graded projective resolution
whose term in homological degree $n$ is generated in total degree $n$.
For a Koszul algebra the quadratic dual is also Koszul, and
\[
 \Hilb_b(A)\Hilb_{-b}(A^!)^{\mathsf T}=\mathrm{Id}_I
\]
by \cite[Proposition 2.9.1 and Lemma 2.11.1]{BeilinsonGinzburgSoergel1996}.

\begin{theorem}\label{thm:chamber-koszul}
For a normal crossing affine oriented matroid, $\Taut$ and $\Inter$
are quadratic and Koszul over every field, and $\Taut^!\cong\Inter$.
\end{theorem}

\begin{remark}
Kowalenko--Mautner ask whether their category $\mathcal O$ algebras
are Koszul for every generic oriented matroid program
\cite[Section 1.3]{KowalenkoMautner}. They prove Koszulity in the
Euclidean case and the numerical Koszul identity in general
\cite[Theorems 7.13--7.14]{KowalenkoMautner}.
Theorem~\ref{thm:chamber-koszul} requires no Euclidean hypothesis,
but concerns the algebras on all feasible chambers. Passing to the
category $\mathcal O$ algebras involves a further choice of chambers
and linear parameter relations, so their Koszulity question requires
an additional argument.
\end{remark}

\begin{corollary}\label{cor:multivariate-inverse}
If $(\M,0)$ is normal crossing, then $V_b^{-1}=\Hilb_{-b}(\Inter)$. Explicitly,
\[
 (V_b^{-1})(P,Q)=(-1)^{d(P,Q)}b_{\Sep(P,Q)}
     \sum_{J\in\Delta_{PQ}}\prod_{i\in J}\frac{b_i^2}{1-b_i^2}.
\]
\end{corollary}

\begin{corollary}\label{thm:varinv}
For a normal crossing affine oriented matroid,
\[
 W_q=(1-q^2)^rV_q^{-1},\qquad
 S_q=(1-q^2)^r(V_q^{-1})^{\cB}.
\]
\end{corollary}

\begin{example}\label{ex:two-points}
For the two-point arrangement of Example~\ref{ex:line},
with its three chambers in order,
\[
 V_b=\begin{pmatrix}
  1&b_1&b_1b_2\\ b_1&1&b_2\\ b_1b_2&b_2&1
 \end{pmatrix},\qquad
 (V_b^{-1})^{\cB}
 =\frac{1-b_1^2b_2^2}{(1-b_1^2)(1-b_2^2)}.
\]
The unique bounded chamber is an interval. Setting $b_1=b_2=q$ gives
\[
 S_q=1+q^2,\qquad
 (V_q^{-1})^{\cB}=\frac{1+q^2}{1-q^2}.
\]
\end{example}

Combining Proposition~\ref{prop:Betti-Varchenko} with
Corollaries~\ref{cor:multivariate-inverse} and~\ref{thm:varinv} gives
an explicit formula for the Betti matrix at $b_0=0$ when $(\M,0)$
is normal crossing.

\begin{corollary}\label{cor:Betti-entries}
If $(\M,0)$ is normal crossing, then for all $P,Q\in\cF$,
\[
 \left.B_b(P,Q)\right|_{b_0=0}
 =(-1)^{r+d(P,Q)}b_{\Sep(P,Q)}
   \sum_{J\in\Delta_{PQ}}\prod_{i\in J}\frac{b_i^2}{1-b_i^2}.
\]
In particular,
\[
 \left.B_b\right|_{b_0=0,\ b_i=q\ (i\in E_0)}
 =(-1)^r(1-q^2)^{-r}W_q.
\]
\end{corollary}

\begin{remark}
If $(\M,0)$ is normal crossing, the Betti matrix at $b_0=0$ is recovered from
the multigraded Hilbert matrix of the intersection algebra $\Inter$:
\[
 \left.B_b\right|_{b_0=0}=(-1)^r\Hilb_{-b}(\Inter).
\]
It would be interesting to extend this categorification beyond normal
crossing, replacing $\Inter$ by the dg Koszul dual of $\Taut$.
We also expect an extension to general $b_0$, giving a categorification
of Lam's full Betti matrix.
\end{remark}


\begin{remark}\label{ex:concurrent-lines}
Consider three distinct lines in the plane meeting at a point.
Their chamber graph is a hexagon. The two taut paths between
opposite chambers have length three and represent the same element
of $\Taut$. The only quadratic relations kill immediate backtracks,
so they cannot identify these two paths. Thus $\Taut$ is not quadratic.

Moreover, for each chamber $P$,
\[
 (W_qV_q)(P,P)=1-2q^2+2q^4-q^6\ne(1-q^2)^2,
\]
so the normal crossing hypothesis in Corollary~\ref{thm:varinv}
cannot be omitted. Lam's face-flag formula still computes $V_b^{-1}$
through his Betti matrix, as recalled in
Proposition~\ref{prop:Betti-Varchenko}.
\end{remark}

\subsection{Proof of quadratic duality}
Let
\[
 \Cube_1=k\bigl(+\mathop{\rightleftarrows}^{u}_{v}-\bigr),\qquad
 \theta=uv+vu,\qquad
 \CubeZero_1=\Cube_1/(\theta)=\Cube_1/(uv,vu).
\]
Identifying the
reversed dual arrows with the arrows having the same endpoints gives
\[
 \Cube_1^!\cong\CubeZero_1,\qquad
 \CubeZero_1^!\cong\Cube_1.
\]
Take one copy for each hyperplane and consider the algebras
\[
 \Cube=\bigotimes_{i\in E_0}\Cube_1,\qquad
 \CubeZero=\bigotimes_{i\in E_0}\CubeZero_1
          =\Cube/(\theta_i\mid i\in E_0),
\]
where $\theta_i$ is the central element $\theta$ in the $i$th factor.
Quadratic duality turns commuting squares into anticommuting
squares, so the dual of an ordinary tensor product is initially a
signed tensor product. Rescaling the arrows removes these signs.
Order $E_0$, write $\sigma^i$ for $\sigma$ with its $i$th sign reversed,
and put
\begin{equation}\label{eq:cube-signs}
 \epsilon_i(\sigma)=(-1)^{|\{j<i\mid\sigma(j)=-\}|}.
\end{equation}
For $i\ne j$,
$\epsilon_i(\sigma)\epsilon_j(\sigma^i)
 =-\epsilon_j(\sigma)\epsilon_i(\sigma^j)$.
Thus mapping the dual arrow $\sigma\to\sigma^i$ to
$\epsilon_i(\sigma)$ times the corresponding arrow of $\Cube$
changes anticommuting squares into commuting squares and gives
\[
 \CubeZero^!\cong\Cube.
\]

\begin{proposition}\label{prop:feasible-models}
Set $e=\sum_{P\in\cF}e_P$. Then
\[
 \Taut\cong e\CubeZero e.
\]
If the arrangement, or affine oriented matroid, is normal crossing,
then also
\[
 \Inter\cong\Cube/\Cube(1-e)\Cube.
\]
\end{proposition}
\begin{proof}
The first identification is left to the reader.
For the second, we compare bases in each block.
Fix $P,Q\in\cF$, and let $\widehat p_{PQ}$ be the taut path class
in $\Cube$. The block $e_P\Cube e_Q$ has basis
\[
 \bigl\{\theta^\nu\widehat p_{PQ}\bigm|\nu\in\NN^{E_0}\bigr\}.
\]
The quotient by $\Cube(1-e)\Cube$ keeps exactly those basis
monomials for which the cube obtained from $P$ by varying the signs
in $\Sep(P,Q)\cup\operatorname{supp}\nu$ has only feasible vertices.

If $(\M,0)$ is normal crossing, Lemma~\ref{lem:feasible-cubes}, applied with
$J=\operatorname{supp}\nu$, identifies this condition with
\[
 \bar P\cap\bar Q\cap H_{\operatorname{supp}\nu}\ne\varnothing,
 \qquad\text{equivalently,}\qquad
 \operatorname{supp}\nu\in\Delta_{PQ}.
\]
The surviving monomials have the same exponent set as the basis
of $R_{PQ}$. Hence the map
$\theta^\nu\widehat p_{PQ}\longmapsto x^\nu 1_{PQ}$ yields the second claimed isomorphism.
\end{proof}

Assume now that $(\M,0)$ is normal crossing. It remains to show that
$e\CubeZero e$ is quadratic.
Every pair of feasible topes is joined by a taut path, so the
arrows between adjacent regions generate $\Taut$, hence the algebra is generated in degree one. To see that quadratic relations suffice, note that any two taut paths with the same endpoints are related by
feasible square moves \cite[Proposition 3.5]{KowalenkoMautner}.
For a non-taut path, square moves expose an immediate backtrack
\cite[proof of Proposition 3.6]{KowalenkoMautner}. These path arguments
use only the feasible topes and the normal crossing condition. Thus commuting
feasible squares and killing immediate backtracks give a quadratic
presentation of $\Taut$.
The corner--quotient lemma \cite[Lemma 8.22]{BradenLicataProudfootWebster2012}
now gives
\[
 \Taut^!
 \cong\CubeZero^!/\CubeZero^!(1-e)\CubeZero^!
 \cong\Cube/\Cube(1-e)\Cube
 \cong\Inter,
\]
proving that the two algebras are quadratic duals.
\subsection{Proof of Koszulity}
Assume that $(\M,0)$ is normal crossing. We adapt the method of the corrected proof
of \cite[Lemma 8.25]{BradenLicataProudfootWebster2012} to all feasible
chambers. We decompose the Koszul complex into multidegrees and use
Lemma~\ref{lem:separating-facets} for the required acyclicity.
Since $\Inter^!\cong\Taut$, it suffices to prove
exactness of the augmented Koszul complex of $\Inter$
\cite[Theorem 2.6.1, \S2.8]{BeilinsonGinzburgSoergel1996}.
Fix $P,R\in\cF$. Multiplying the Koszul complex by the corresponding
idempotents gives
\[
 K_n=e_P\Inter\otimes_{\Inter_0}(\Taut_n)^*e_R
 =\bigoplus_{\substack{Q\in\cF\\d(Q,R)=n}}
     e_P\Inter e_Q\otimes_k k\tau_{RQ}^*.
\]
The differential preserves multidegree, so we can further fix
$\mathbf m\in\NN^{E_0}$ and work in the subcomplex
$K_{\bullet,\mathbf m}$.

To prove exactness, we distinguish two cases. If some $m_i\geq2$,
one component of the differential pairs the basis vectors.
If all $m_i\leq1$, we identify the nonzero component with the
reduced cochain complex of $L_P^R$, defined before
Lemma~\ref{lem:separating-facets}.

Using the signs in \eqref{eq:cube-signs}, we decompose the differential as
$\partial=\sum_{i\in E_0}\partial_i$, where
\[
 \partial_i(z\otimes\tau_{RQ}^*)=
 \begin{cases}
  \epsilon_i(Q)\,z 1_{Q,Q^i}\otimes\tau_{R,Q^i}^*,
     &i\in\Sep(Q,R),\ Q^i\in\cF,\\
  0,&\text{otherwise}.
 \end{cases}
\]
Moreover, we will denote
\begin{equation}\label{eq:koszul-exponent}
 \nu(Q)=\tfrac12\bigl(\mathbf m-\one_{\Sep(P,Q)}
                              -\one_{\Sep(Q,R)}\bigr),
\end{equation}
so that $K_{n,\mathbf m}$ has basis vectors
\[
 v_Q=x^{\nu(Q)}1_{PQ}\otimes\tau_{RQ}^*,\qquad
 \begin{cases}
  Q\in\cF,\quad d(Q,R)=n,\\
  \nu(Q)\in\NN^{E_0},\\
  \operatorname{supp}\nu(Q)\in\Delta_{PQ}.
 \end{cases}
\]
The two indicator entries in \eqref{eq:koszul-exponent} sum to
$1$ where $P(i)\ne R(i)$, and to $0$ or $2$ otherwise.
The component $K_{n,\mathbf m}$ is zero unless
$\mathbf m\equiv\one_{\Sep(P,R)}\pmod2$. Hence we assume this parity condition from now on.
For $i\in\Sep(Q,R)$, \eqref{eq:koszul-exponent} and
\eqref{eq:intersection-product} give
\[
 \nu(Q^i)-\nu(Q)=
 \begin{cases}
  \mathbf e_i,&P(i)=R(i),\\
  0,&P(i)\ne R(i),
 \end{cases}
 \qquad
 \partial_i v_Q=
 \begin{cases}
  \epsilon_i(Q)v_{Q^i},&\text{if }v_{Q^i}\text{ exists},\\
  0,&\text{otherwise}.
 \end{cases}
\]

\smallskip
\noindent\emph{Case 1: $m_i\geq2$ for some $i$.}
The basis vectors occur in pairs $v_Q,v_{Q^i}$, with
$\partial_i v_Q=\epsilon_i(Q)v_{Q^i}$ when $Q(i)\ne R(i)$.
Since $m_i\geq2$, formula~\eqref{eq:koszul-exponent} gives
$\nu(Q^i)\in\NN^{E_0}$ for every basis vector $v_Q$, and
\[
 \Sep(P,Q)\cup\operatorname{supp}\nu(Q)
 =\Sep(P,Q^i)\cup\operatorname{supp}\nu(Q^i).
\]
By Lemma~\ref{lem:feasible-cubes}, the cube based at $P$ with this
common set of directions is feasible. It contains $Q^i$, and the
same lemma gives $\operatorname{supp}\nu(Q^i)\in\Delta_{P,Q^i}$.
Hence, every partner $v_{Q^i}$ exists.
In each homological degree, set
\[
 U_n=\langle v_Q\mid Q(i)\ne R(i)\rangle_k,\qquad
 V_n=\langle v_Q\mid Q(i)=R(i)\rangle_k.
\]
The pairing gives $\partial_i:U_n\xrightarrow{\sim}V_{n-1}$.
Every $\partial_j$ with $j\ne i$ preserves $Q(i)$, hence $U$ and $V$.
For a cycle $z=u+v\in U_n\oplus V_n$, choose
$y\in U_{n+1}$ with $\partial_i y=v$.
Then $w=z-\partial y$ is a cycle in $U_n$.
The $V_{n-1}$-component of $\partial w$ is $\partial_i w$, so
injectivity of $\partial_i$ gives $w=0$ and $z=\partial y$.

\smallskip
\noindent\emph{Case 2: every $m_i\leq1$.}
The parity condition gives
$\mathbf m=\one_{\Sep(P,R)}$.
A basis vector $v_Q$ now has $\nu(Q)=0$ and
\[
 \Sep(P,Q)\uplus\Sep(Q,R)=\Sep(P,R).
\]
Let $G$ be the covector of the nonempty
common face $\bar P\cap\bar Q$. Corollary~\ref{cor:face-subsets}
gives $z(G)=\Sep(P,Q)$.
Associate $v_Q$ with the simplex $z(G)$ of $L_P^R$.
The inverse recovers $Q=G\circ R$, a feasible tope with
$G\leq P,Q$ and $\Sep(P,Q)=z(G)$.
The basis of $K_{n,\mathbf m}$ therefore corresponds to the simplices
with $|z(G)|=d(P,R)-n$.

A nonzero differential term in direction $i$ sets $G(i)$ to zero,
so it adds $i$ to the simplex $z(G)$.
Rescaling $v_Q$ by $\prod_{j\in z(G)}\epsilon_j(P)$
gives the usual simplicial coboundary signs. This component is the reduced
cochain complex
\[
 K_{n,\mathbf m}
 \cong\widetilde C^{\,d(P,R)-n-1}(L_P^R;k).
\]
For $P\ne R$, Lemma~\ref{lem:separating-facets} makes it exact.
For $P=R$, only $v_P$ remains, and the augmentation
maps it isomorphically to $k$.
